\lysh{33891}{4}

\begin{alist}
\item Δίνεται η γεωμετρική πρόοδος $(\alpha_\nu)$ με λόγο $\lambda$ για την οποία ισχύουν:
$\alpha_3 = 4, \alpha_5 = 16$ και $\lambda > 0$.
Έχουμε:
\[\begin{cases} \alpha_3 = 4 \\ \alpha_5 = 16 \end{cases} \Leftrightarrow \begin{cases} \alpha_1 \cdot \lambda^2 = 4 \\ \alpha_1 \cdot \lambda^4 = 16 \end{cases} \stackrel{(:)}{\Leftrightarrow} \begin{cases} \lambda^2 = 4 \\ \alpha_1 \cdot \lambda^2 = 4 \end{cases} \stackrel{\lambda > 0}{\Leftrightarrow} \begin{cases} \lambda = 2 \\ \alpha_1 \cdot 2^2 = 4 \end{cases} \Leftrightarrow \begin{cases} \lambda = 2 \\ \alpha_1 = 1 \end{cases}\]

\item Έχουμε
$\dfrac{\beta_{\nu+1}}{\beta_\nu} = \dfrac{\dfrac{1}{\alpha_{\nu+1}}}{\dfrac{1}{\alpha_\nu}} = \dfrac{\alpha_\nu}{\alpha_{\nu+1}} = \dfrac{1}{\lambda}$, οπότε $(\beta_\nu)$, με $\beta_\nu = \dfrac{1}{\alpha_\nu}, \nu = 1, 2, 3, \dots$ είναι επίσης
γεωμετρική πρόοδος με πρώτο όρο $\beta_1 = \dfrac{1}{\alpha_1} = 1$ και λόγο τον αντίστροφο του λόγου της
$(\alpha_\nu)$, δηλαδή $\lambda' = \dfrac{1}{2}$.

\item Εφόσον $S_{10}$ είναι το άθροισμα των $10$ πρώτων όρων της $(\alpha_\nu)$, έχουμε:
\[S_{10} = \alpha_1 \dfrac{\lambda^{10} - 1}{\lambda - 1} = 1 \cdot \dfrac{2^{10} - 1}{2 - 1} = 2^{10} - 1.\]
Το άθροισμα των $10$ πρώτων όρων της $(\beta_\nu)$ είναι:
\[S'_{10} = \beta_1 \dfrac{(\lambda')^{10} - 1}{\lambda' - 1} = 1 \cdot \dfrac{\left(\dfrac{1}{2}\right)^{10} - 1}{\dfrac{1}{2} - 1} = \dfrac{\dfrac{1 - 2^{10}}{2^{10}}}{-\dfrac{1}{2}} = \dfrac{2^{10} - 1}{2^9} = \dfrac{S_{10}}{2^9}.\]
\end{alist}

